\documentclass{article}
\usepackage{tmlr}
% \usepackage[accepted]{tmlr}   % camera-ready
\usepackage{amsmath,amssymb,amsthm}
\usepackage{booktabs}
\usepackage{graphicx}
\usepackage{xcolor}
\usepackage{hyperref}
\usepackage{url}
\usepackage{multirow}

\newtheorem{proposition}{Proposition}
\newtheorem{remark}{Remark}
\newcommand{\eps}{\varepsilon}
\newcommand{\ACT}{\mathrm{ACT}}
\newcommand{\Jfifty}{J_{50}}

\title{Certifying Retrieval Policies with Non-Neutral AI Judges:\\ What a Judge Saves, and How a Pilot Predicts It}

\author{\name Anonymous authors \\ \addr Paper under double-blind review}

\begin{document}
\maketitle

\begin{abstract}
Before a retrieval policy is deployed in a new domain, an auditor must certify from a limited human audit that it is within a tolerance $\eps$ of the best alternative. LLM relevance judges can label every candidate document, but may be biased toward the policies they judge. We carry prediction-powered and active inference to this certification problem and measure what they deliver on eight retrieval collections, five AI judges and three pre-registered evaluations. The value of a judge is set by the correlation $\rho$ between true and judge-predicted paired policy differences, not by its label accuracy: over 54 settings the realised gain correlates $0.75$ with $\rho$ and $0.03$ with accuracy. A sample-split certificate uses the judge only as a mean-zero control variate, so fixed policy contrasts are estimated without bias for any judge; coverage rests on an asymptotically calibrated $t$ bound whose observed wrong-certificate rate is at most $0.013$. With every human label charged once, pilot included, importance sampling by closed-form per-document decision weights needs $40$--$56\%$ fewer labels than uniform sampling; a $\lambda$-fitted judge control variate saves a further $3$--$20\%$ where $\rho$ is high and nothing otherwise; an oracle allocation across the comparisons of a menu saves $0$--$14\%$ at $\eps=0.02$. A single cost model, saving $=$ variance reduction $\times$ post-pilot share of the cost, fits 193 cells with correlation $0.95$, and both factors can be computed from a 20-query pilot. In tests with predictions fixed before the audit, including one collection never used in development, the pilot predicted the relative cost of designs (correlation $0.83$--$0.97$) but over-estimated the individual judge's gain, and a single-pilot adoption rule missed the usefulness criteria set for it. What a judge saves is thus measurable and, in its relative part, pilot-predictable; whether a given judge pays on a new collection is what a 20-query pilot answers least reliably.
\end{abstract}

\section{Introduction}
Retrieval systems are increasingly deployed as configurable policies: a fixed ranker plus a rule that decides how many documents to return, which threshold to apply, or which reranker to trust. When such a system moves to a new domain, an auditor with a small human-judgment budget must decide whether a candidate policy can be deployed, that is, whether its utility is within $\eps$ of the best alternative on the target distribution, and otherwise collect more judgments or abstain. LLM relevance judges can now label every candidate document, so the audit can be assisted by a judge. The policies, however, are often built from the same model families as the judges, and a judge need not be neutral toward what it evaluates \citep{balog2025rankers}. Prediction-powered inference (PPI) removes the bias of an AI predictor with human labels, and active statistical inference tells where to spend those labels for a given functional \citep{angelopoulos2023ppi,angelopoulos2023ppipp,zrnic2024active}. What they do not tell an auditor is which judge is worth using for a policy decision, what the simultaneous certification of a menu of policies will cost in human labels once the pilot is charged, and whether that cost can be foreseen from the pilot. This paper carries these tools to policy certification, measures what they deliver, and reports the outcome of testing the answers on data locked before analysis.

\paragraph{Contributions.}
\begin{enumerate}
\item \textbf{What a judge is worth is not its accuracy.} The PPI++ gain for a paired policy difference is set by the correlation $\rho$ between true and judge-predicted differences and by the labeled-to-unlabeled ratio (Proposition~\ref{prop:A}, a restatement of PPI++). Across 54 (collection, pair, judge) settings the realised gain correlates $0.75$ with $\rho$ and $0.03$ with judge accuracy (Section~\ref{sec:judge}); a mechanism map shows where a certificate gains, and a simulation reproduces it from $\rho$, the gap, $n$, $N$ and $\eps$. A judge built from the model that defines a policy over-credits it by $0.14$ F1; we separate this mean bias, which the correction removes, from the loss of $\rho$, which costs efficiency (Section~\ref{sec:selfpref}).
\item \textbf{A certificate that keeps PPI's judge-independence under data-driven selection.} A sample-split certificate uses the judge only as a mean-zero control variate, as PPI does; the estimator of every fixed policy contrast is therefore unbiased for any judge, and a data-dependent candidate is covered by simultaneous bounds over the fixed family of contrasts (Proposition~\ref{prop:B}). Its coverage rests on the bound, an asymptotically calibrated $t$ bound whose observed wrong-certificate rate is $\le0.008$ on the development collections and $\le0.013$ in every later test; an exact betting certificate is affordable on two collections and close to a census on two others, because of its range term (Section~\ref{sec:valid}).
\item \textbf{Where to spend human judgments, and a cost model that predicts it.} The paired difference is a linear functional with closed-form per-document decision weights (Proposition~\ref{prop:C}). In the document-level, precision-at-cutoff experiments, with every unique label charged once, pilot included, importance sampling by these weights needs $40$--$56\%$ fewer labels than uniform sampling; a $\lambda$-fitted judge control variate saves a further $3$--$20\%$ of the total where its $\rho$ is high and nothing otherwise; an oracle allocation across the comparisons of a menu saves only $0$--$14\%$ at $\eps=0.02$ ($1$--$17\%$ at $\eps=0.01$). A single cost model accounts for these numbers, a design's saving being its variance reduction times the post-pilot share of the cost; it fits 193 cells with correlation $0.95$, and both factors can be computed from the 20-query pilot. Tested with predictions fixed in advance, the pilot predicts the relative cost of designs on three further collections; on the one never used in development it over-estimates the judge's own gain, and the single-pilot adoption rule misses the usefulness criteria set for it (Section~\ref{sec:where}).
\item \textbf{Three pre-registered external evaluations.} Validity replicated in all three; the pre-registered efficiency criteria were not met in any, and Section~\ref{sec:prereg} reports why.
\end{enumerate}
We propose no new estimator or sampling optimality result; the estimator is a special case of active inference for a linear functional (Section~\ref{sec:relation}). Nor do we offer a general theory of judge neutrality.

\section{Setup}\label{sec:setup}
Let $P$ be the target query distribution, $\mathcal{M}$ a finite menu of policies acting on a shared ranked candidate pool for each query $q$, and $u_m(q)\in[0,1]$ the utility of policy $m$ on $q$ given human relevance labels $y_d\in\{0,1\}$ for the pooled documents $d$. Two utilities are used: set-F1, $u_m(q)=2\,\mathrm{tp}_m/(k_m+n_G)$ with $k_m=|R_m(q)|$ and $n_G(q)=\sum_d y_d$, and precision at the policy's cutoff, $u_m(q)=\mathrm{tp}_m/k_m$. Population utilities are $\mu_m=\mathbb{E}_P[u_m(q)]$.

\paragraph{Decision and error.} A certificate $\ACT$ selects a candidate $\hat m$ and asserts $\max_{j}\mu_j-\mu_{\hat m}\le\eps$; a \emph{wrong certificate} is an $\ACT$ whose true regret exceeds $\eps$. For looks $t=1..L$ and ordered pairs $(j,\hat m)$, each one-sided upper confidence bound $U_{t,j}$ on $\mu_j-\mu_{\hat m}$ is at level $1-\alpha'$ with $\alpha'=\alpha/(L\,|\mathcal M|(|\mathcal M|-1))$, and $\ACT_t\iff\max_{j\ne\hat m}U_{t,j}\le\eps$ ($L=1$ in the fixed-budget designs of Section~\ref{sec:where}). Throughout $\alpha=0.10$ and $\eps\in\{0.005,0.01,0.02,0.05\}$.

\paragraph{Audit units and the judge.} A \emph{query-level} audit labels a query's whole pool (needed for set-F1 because of $n_G$); a \emph{document-level} audit labels a sampled subset of documents. An AI judge is a fixed map $f$ from a (query, document) pair to a relevance probability, thresholded at $0.5$ for judge utilities $\hat u_m(q)$. Nothing below assumes the judge is calibrated, unbiased, or independent of the policies.

\paragraph{Efficiency metric.} In the \emph{sequential} query-level design (Sections~\ref{sec:valid} and \ref{sec:prereg}) one audit path is followed through the looks and we report the cumulative $\ACT$ rate. In the \emph{fixed-budget} document-level design (Section~\ref{sec:where}) each budget is a separate audit with its own pilot and sample, evaluated once ($L=1$); $\Jfifty$ is the budget, in unique human-judged (query, document) pairs with every label consumed included, at which the certification probability first reaches 50\%, interpolated between the budgets run. $\Jfifty$ is the cost of a pre-committed budget that certifies half the time, not the stopping cost of a sequential procedure. A cell whose largest budget stays below 50\% is \emph{not reached} and never enters a ratio. Query-level (set-F1) and document-level (precision) results use different utilities and are never compared to each other.

\section{What determines the value of a judge}\label{sec:judge}
\begin{proposition}[PPI gain for paired differences]\label{prop:A}
Fix policies $j,m$ and let $D=u_j-u_m$, $\hat D=\hat u_j-\hat u_m$ with variances $\sigma^2,\hat\sigma^2$ and correlation $\rho$. With $n$ audited and $N$ unlabeled queries drawn independently, the PPI++ estimator $\hat\mu(\lambda)=\lambda\,\overline{\hat D}_N+\overline{(D-\lambda\hat D)}_n$ has variance $[\sigma^2-2\lambda\rho\sigma\hat\sigma+\lambda^2\hat\sigma^2]/n+\lambda^2\hat\sigma^2/N$, minimised over unconstrained $\lambda$ at $\lambda^\star=\rho\sigma/\big(\hat\sigma\,(1+n/N)\big)$ to $\frac{\sigma^2}{n}\big(1-\frac{\rho^2}{1+n/N}\big)$. The effective-sample-size gain over the human-only estimator is $G=1/\big(1-\rho^2/(1+n/N)\big)$, which increases to $1/(1-\rho^2)$ as $N/n\to\infty$.
\end{proposition}
This is PPI++ \citep{angelopoulos2023ppipp} applied to a paired difference (proof in Appendix~\ref{app:proofs}). We use it as the reference curve of Figure~\ref{fig:F4}, with three caveats. $G$ is the gain at the oracle $\lambda^\star$, whereas the implementation cross-fits $\lambda$ on halves of the audited queries and loses part of the gain to estimation noise at small $n$. The implementation clips $\lambda$ to $[0,1]$, so a negatively correlated judge is switched off rather than exploited. Finally, $G$ is a variance ratio, and a certificate is a threshold event on a bound (Figure~\ref{fig:F6}).

\paragraph{Bias toward a policy and $\rho$ are different quantities.} Writing $\hat u_m=u_m+b_m+e_m$ with a policy-specific bias $b_m$ and a query-level error $e_m$, $\hat D=D+(b_j-b_m)+(e_j-e_m)$. A constant $b_j-b_m$, however large, leaves $\rho$ unchanged ($\hat D=D+0.1$ has $\rho=1$) and is removed by the control variate of Section~\ref{sec:valid}; what lowers $\rho$ is the query-varying part of the error relative to $\mathrm{Var}(D)$. Bias toward a policy therefore has two consequences that should not be conflated. A \emph{mean bias} corrupts judge-only comparisons but not the corrected estimator; a \emph{loss of predictive correlation} costs efficiency. The second need not accompany the first (Section~\ref{sec:selfpref}).

\paragraph{Empirical test.} We measured $\rho$ and the realised gain, defined as the ratio of squared estimated standard errors of the human-only and PPI++ estimators over 300 random audits of $T=10$ queries with cross-fitted $\lambda$, on 54 (collection, pair, judge) settings: four BEIR development pools under two retrieval stacks, five fully judged pools, and two judges (Qwen3-Reranker-0.6B, Qwen3-8B). Figure~\ref{fig:F4} shows $\mathrm{corr}(\text{gain},\rho)=0.75$ against $\mathrm{corr}(\text{gain},\text{pair accuracy})=0.03$. Demeaned within the 14 (collection, judge) groups the correlation is $0.66$ and positive in 13 of 14 (the exception has every $\rho$ between $-0.17$ and $0.07$); in all 13 settings where switching from the reranker to the LLM judge raised $\rho$, the gain rose. On the 33 settings whose pools are released, the Monte-Carlo MSE ratio about the population gap correlates $0.78$ with $\rho$ and agrees with the estimated-SE ratio to a median $0.03$ (Appendix~\ref{app:repro}).

\begin{figure}[t]\centering
\includegraphics[width=0.72\linewidth]{figures/F4_gain_vs_rho.png}
\caption{Realised PPI gain (ratio of squared estimated standard errors) at $T=10$ against $\rho$, 54 settings, cross-fitted $\lambda$. The dashed curve is $1/(1-\rho^2)$. Gains track $\rho$ (0.75) and not accuracy ($0.03$).}\label{fig:F4}
\end{figure}

\paragraph{When does the gain materialise?} On DBpedia-Entity (399 queries) and TREC DL 2021--23 (211 queries) we varied the judge's $\rho$ (replacing a fraction of its labels with base-rate coin flips: $\rho=0.64,0.57,0.39,0.26$ on DBpedia-Entity), the unlabeled size $N_{\mathrm{unlab}}\in\{50,100,200,300\}$ and $\eps$, and measured the $\ACT$ ratio of the PPI certificate to the human-only certificate at $T=90$ (Figure~\ref{fig:F6}; 300 repeats per cell). With $\rho\le0.39$ the ratio lies between $0.76$ and $1.06$ in every cell. With $\rho\ge0.57$ and $\eps\le0.02$ it exceeds $1.2$ in 11 of the 12 cells with $N_{\mathrm{unlab}}\ge200$ ($1.21$--$1.53$) and, for $\rho=0.64$, already at $N_{\mathrm{unlab}}=100$ and $\eps=0.01$ ($1.38$); at $N_{\mathrm{unlab}}=50$ it never exceeds $1.17$, and at $\eps=0.05$, where humans alone certify about 90\% of runs, it is $0.95$--$1.01$. Every cell has wrong-certificate rate $\le0.01$. We read the map as a description of where the gain appeared under these conditions rather than as a universal threshold; $\rho$ matters first, the unlabeled pool second, and a tolerance that humans cannot settle alone is a precondition. On TREC DL 2021--23 the judge reaches $\rho\le0.46$ with $N_{\mathrm{unlab}}\le100$, and no cell exceeds $1.05$, in line with the pre-registered evaluation of Section~\ref{sec:prereg}. The map can be reproduced from a few population quantities. Running the same certificate code on Gaussian differences with the population's gaps, standard deviations and $\rho$ gives the 64 DBpedia-Entity cells with correlation $0.85$ (agreement on the $1.2$ threshold in 92\% of cells) and the 32 TREC DL cells with mean absolute error $0.03$. The reproduction requires the cross-fitted $\lambda$; with the population-optimal $\lambda$ the poor judge's losses on TREC DL disappear, so those losses are due to estimating $\lambda$ on 22-query folds (Appendix~\ref{app:calculator}).

\begin{figure}[t]\centering
\includegraphics[width=\linewidth]{figures/F6_rho_map.png}
\caption{Mechanism map: $\ACT(\text{PPI})/\ACT(\text{humans})$ at $T=90$ over $\rho$ (judge-label noise), unlabeled pool size and $\eps$; wrong-certificate rate $\le0.01$ in every cell. Gains appear only for $\rho\ge0.57$ and $\eps\le0.02$, grow with $N_{\mathrm{unlab}}$, and are absent at $N_{\mathrm{unlab}}=50$ and at $\eps=0.05$.}\label{fig:F6}
\end{figure}

\subsection{Judges built from the model that defines a policy}\label{sec:selfpref}
We added a fourth policy, \texttt{rr\_thresh}, which returns every pooled document whose Qwen3-Reranker score exceeds a tuned threshold, so that one policy is defined by a model in the judges' family, and compared five judges on the six policy pairs of five fully judged pools at fixed policies (Table~\ref{tab:selfpref}). Two signatures are specific to the \emph{identical-model} judge: it over-credits its own policy by $0.14$ on the paired F1 difference (other judges $-0.04$ to $+0.03$), and its error rate on the documents where the two policies disagree rises from $0.36$ to $0.41$. The drop in $\rho$ for pairs involving \texttt{rr\_thresh} appears for every judge, including the different-family MPNet rule, so it is a property of that policy's paired differences rather than of self-preference. Validity is unaffected ($0/500$ wrong certificates with the four-policy menu for both Qwen judges); the identical-model judge yields no efficiency gain (ACT $246$ vs $245$ human-only) while the 8B judge does ($328$).

\begin{table}[t]\centering\footnotesize
\caption{Self-preference test at fixed policies: mean over 15 policy pairs per judge, split by whether the pair involves the reranker-defined policy (``without / with''). ``Bias'' is the mean of predicted minus true paired difference.}\label{tab:selfpref}
\begin{tabular}{lcccc}\toprule
Judge & $\rho$ & in-band / out-of-band error & bias & PPI gain\\\midrule
Qwen3-Reranker (identical model) & 0.57 / 0.25 & 0.36 / 0.32 $\to$ 0.41 / 0.31 & $-0.03$ / $\mathbf{-0.14}$ & 1.07 / 0.95\\
Qwen3-8B (same family) & 0.69 / 0.48 & 0.28 / 0.24 $\to$ 0.29 / 0.23 & $0.00$ / $-0.04$ & 1.19 / 1.04\\
Mistral-7B (other family) & 0.38 / 0.30 & 0.34 / 0.34 $\to$ 0.37 / 0.33 & $+0.02$ / $-0.04$ & 1.10 / 1.06\\
MPNet rank rule (other family) & 0.39 / 0.21 & 0.39 / 0.35 $\to$ 0.40 / 0.35 & $-0.05$ / $+0.03$ & 1.03 / 0.94\\\bottomrule
\end{tabular}
\end{table}

\section{A certificate whose bias does not depend on the judge}\label{sec:valid}
Our earlier certificate chose the candidate on the data used for the bound (leave-one-out cross-fitted utilities, paired bootstrap, Bonferroni over looks $\times(|\mathcal M|-1)$ comparisons). On synthetic data this alone raises the per-look miss rate from $0.020$ to $0.049$. On two development pools its wrong-certificate rate was $0.146$ $[0.116,0.180]$ and $0.180$ $[0.147,0.217]$ at $\alpha=0.10$.

\paragraph{Split certificate.} At look $t$ the audited queries are split into a training half $A_t$ and a validation half $B_t$. Policy parameters and the judge-adoption decision are made on $A_t$; the candidate $\hat m_t$ is the policy with the largest lower confidence bound on $B_t$; the bound is computed on $B_t$ only, as a one-sided $t$ bound on the paired differences (humans only) or as a PPI++ bound with $\lambda$ cross-fitted inside $B_t$ and the unlabeled mean taken outside $B_t$ (with the judge).
\begin{proposition}[Simultaneous validity, judge-independent bias]\label{prop:B}
Fix the menu $\mathcal M$ before the audit. Suppose that, for every look $t$ and every ordered pair $(j,m)$, $U_{t,j,m}$ computed on $B_t$ is a level-$(1-\alpha')$ upper confidence bound for $\mu_j-\mu_m$ with $\alpha'=\alpha/(L|\mathcal M|(|\mathcal M|-1))$. Then on the event $\mathcal E=\bigcap_{t,j\ne m}\{\mu_j-\mu_m\le U_{t,j,m}\}$, which has probability at least $1-\alpha$, every data-dependent candidate $\hat m_t$---chosen on $A_t$, on $B_t$, or on both---satisfies $\mu_j-\mu_{\hat m_t}\le U_{t,j,\hat m_t}$ for all $j$; hence $\Pr(\exists t:\ \ACT_t\wedge\mu_{\mathrm{best}}-\mu_{\hat m_t}>\eps)\le\alpha$. The statement holds for every fixed judge $f$, because for every \emph{fixed} ordered pair $(j,m)$ the estimator behind $U_{t,j,m}$ is unbiased for $\mu_j-\mu_m$ conditionally on the training half, the judge entering only through a control variate whose expectation is zero; the candidate enters only through which member of the covered family is read off.
\end{proposition}
The proof (Appendix~\ref{app:proofs}) is a union bound over ordered pairs and looks, together with the independence of a cross-fitted $\lambda$ from the fold it is applied to. The proposition contains three claims with different hypotheses. Unbiasedness of every fixed contrast needs only the split and no property of the judge; the estimator of the \emph{selected} contrast is not claimed to be unbiased and need not be. Safe certification of a candidate chosen on any part of the data needs simultaneous coverage of the fixed family. An error of at most $\alpha$ in finite samples needs each bound to hold its level in finite samples, which the $t$ bound does only asymptotically: a difference equal to $1$ with probability $0.02$ gives a $t$ bound of $0$ whenever no success is seen, which happens on $55\%$ of samples at $n=30$. The certificate is therefore asymptotically calibrated, and the wrong-certificate rates below are observed Monte-Carlo rates rather than guarantees. The adoption rule never enters the bound; a poor adoption decision changes which valid bound is used and therefore only efficiency.

\paragraph{Empirical validity.} Across 24 (collection, judge, utility) development cells at 500 repeats, the split certificate's wrong-certificate rate is $\le 4/500=0.008$. With an \emph{adversarial} judge (inverted reranker, pair accuracy $0.34$) it is $2/500$ and the PPI certificate degenerates to the human-only one ($156$ vs $158$ cumulative $\ACT$ at $T=190$) because $\lambda\to0$ (Figure~\ref{fig:F5}, Appendix~\ref{app:exact}). Two stress designs for the document-level arms of Section~\ref{sec:where}, one forcing the candidate to the true worst policy and a \emph{boundary} design with the runner-up as candidate and $\eps$ at $0.9\times$ its true regret so that every $\ACT$ is an error, give type-I error $\le0.03$ for every arm.

\paragraph{What exactness costs.} Fixing the estimand to the mean over the $N$ queries of the evaluation set (the training half is then known exactly and the validation half is a sample without replacement), we compared at $\eps=0.01$ on the four target collections ($N=159$--$399$) the pre-registered $t$ bound, the same bound with the finite-population correction, and the exact without-replacement betting bound of \citet{waudby2020wor} (Appendix~\ref{app:exact}). The finite-population $t$ certificate stops $20$--$27\%$ earlier than the pre-registered one at no observed cost in error. The betting certificate fires eventually in $88$--$100\%$ of runs, but on DBpedia-Entity and DL 21--23 it has certified in $90\%$ and $73\%$ of runs by the time $75\%$ of the queries are audited, whereas on ANTIQUE and CAsT it fires almost only at the last look and saves essentially nothing over a census. The cause is the range term of the bound, not its variance: with the loose range $[-1,1]$ and $\alpha'\approx10^{-3}$, the second-order term $R\log(1/\alpha')/n$ alone exceeds $\eps$ until $n\gtrsim1{,}400$. We therefore keep the $t$ certificate for the main results and report its error rate as observed.

\paragraph{Always using the judge can cost efficiency; a pilot rule protects it.} On TREC DL 2021--23 with the reranker judge (pair accuracy $0.43$), always-PPI reached $256/500$ cumulative $\ACT$ at $T=90$ against $313/500$ for humans only; the finite-$N$ optimal $\lambda$ recovered $282$ and a variance guard $266$. A rule that adopts the judge only when a BCa lower bound on $\rho$ from the training half (at least 30 queries) implies a gain above $1.1$ adopted the reranker in $1.3\%$ of runs and stayed at $309$. The rule protects efficiency and saves judge inference. Validity does not depend on it, and its bound is an adoption score rather than a guarantee.

\section{Where to spend human judgments}\label{sec:where}
\begin{proposition}[Decision weights of documents]\label{prop:C}
For prefix policies $a,b$ with cutoffs $k_a\le k_b$ on the same ranked pool and a utility $u_m=\sum_{d\in R_m}y_d/Z_m$ whose normaliser does not depend on the labels, $u_a-u_b=\sum_d w(d)\,y_d$ with $w(d)=1/Z_a-1/Z_b$ on common documents, $w(d)=-1/Z_b$ on the band $R_b\setminus R_a$, and $0$ elsewhere. Exact band sufficiency holds iff $Z_a=Z_b$; otherwise the common documents carry weight $|1/Z_a-1/Z_b|$, small when the cutoffs are close and never zero. For set-F1 the normaliser depends on the labels through $n_G$ and the relation holds only after linearisation (Appendix~\ref{app:f1}).
\end{proposition}
Band sufficiency does not hold for precision: with $R_A=\{a\}$, $R_B=\{a,b\}$ and $b$ irrelevant, the difference is $0.5$ or $0$ depending on $a$ alone. A band-only audit built on the opposite assumption gave no gain in our runs (Section~\ref{sec:negative}). With $\pi_d$ the inclusion probability of a human label and $\hat\jmath_d$ the judge's label, the Horvitz--Thompson estimator $\hat D_{\mathrm{HT}}=\sum_{d\text{ sampled}}w(d)y_d/\pi_d$ and the control-variate estimator $\hat D_{\mathrm{CV}}=\sum_d w(d)\hat\jmath_d+\sum_{d\text{ sampled}}w(d)(y_d-\hat\jmath_d)/\pi_d$ are unbiased for $D(q)$ for any judge, and under Poisson sampling (each document included independently with probability $\pi_d$, which is what the implementation does) their variances are $\sum_d w(d)^2y_d^2(1-\pi_d)/\pi_d$ and $\sum_d w(d)^2(y_d-\hat\jmath_d)^2(1-\pi_d)/\pi_d$; other without-replacement designs add covariance terms. The judge helps on the documents where it is right, weighted by decision weight, and never affects unbiasedness. Query-level certificates (Proposition~\ref{prop:B}) then apply with $\hat D$ in place of a fully audited $D$.

\subsection{Relation to active statistical inference}\label{sec:relation}
$\hat D_{\mathrm{CV}}$ is the active (prediction-powered) estimator of \citet{zrnic2024active} written for the linear functional $\sum_d w(d)y_d$. Under a budget $\sum_d\pi_d=b$ its variance-minimising allocation is $\pi_d^\star\propto|w(d)|\sqrt{\mathbb E[(y_d-\hat\jmath_d)^2]}$. Our rule $\pi_d\propto|w(d)|$ is the special case of a constant residual model; the ``calibrated-judge'' rule uses $\hat\jmath_d(1-\hat\jmath_d)$; the ``residual-estimated'' rule fits $\mathbb E[(y-\hat\jmath)^2]$ on the pilot; \citet{li2025robust} mix an uncertainty rule with a safe distribution, which here is $\pi_d\propto|w(d)|$ itself. What is specific to this paper is the functional (the decision weights) and the certificate that consumes it.

\subsection{Comparison at equal human budget}\label{sec:table2}
All arms share a 20-query fully labelled pilot that fixes policy parameters, the candidate and any residual model, the same expected number of labelled documents per round, and the same certificate. Every human label is charged once through a ledger keyed by (query, document): the pilot is charged for every document of its 20 queries, and one sampled document serves all comparisons of the menu (the per-comparison accounting of the pre-registered runs is kept in Appendix~\ref{app:legacy} as the historical record). The certificate and the wrong-certificate score target the same estimand, the mean paired difference over all $N$ queries of the collection: the pilot's contribution is known exactly and the remaining queries are a without-replacement sample whose document-level estimates carry a two-stage variance (Appendix~\ref{app:costmodel}). Each budget is a separate audit with a fresh pilot and sample, evaluated once, so Table~\ref{tab:J50} prices pre-committed budgets, not a sequential procedure. It reports $\Jfifty$ for precision at cutoff at $\eps=0.02$ on the four fully judged target collections; at $\eps=0.01$ (Table~\ref{tab:J50b}) every arm except uniform sampling on CAsT and DBpedia-Entity reaches 50\% $\ACT$, with the same ordering.

\begin{table}[t]\centering\scriptsize\setlength{\tabcolsep}{2.2pt}
\caption{$\Jfifty$: \emph{unique} human-labelled (query, document) pairs, pilot included, at which the certification probability of a pre-committed budget reaches 50\% (precision at cutoff, $\eps=0.02$, $\alpha=0.10$, 300 draws). Observed wrong-certificate rate $\le0.003$ in every cell; no cell is censored. Judges: Qwen3-8B (Q), Qwen3-Reranker-0.6B (R), inverted reranker (adversarial, I). The two bold rows are the human-only reference arm and the judge arm the text discusses. Monte-Carlo standard error of a cell is about 6--9\%; paired-bootstrap 95\% intervals of the savings ratios are given in the text.}\label{tab:J50}
\begin{tabular}{l rrr rrr rr r}\toprule
 & \multicolumn{3}{c}{ANTIQUE} & \multicolumn{3}{c}{CAsT 2019} & \multicolumn{2}{c}{DBpedia-Entity} & DL 21--23\\
Method & Q & R & I & Q & R & I & Q & R & Q\\\midrule
Uniform document sampling (humans) & 1,405 & 1,405 & 1,405 & 4,778 & 4,778 & 4,778 & 5,563 & 5,563 & 4,092\\
Uniform over decision-relevant documents (humans) & 1,107 & 1,107 & 1,107 & 2,823 & 2,823 & 2,823 & 3,372 & 3,372 & 2,613\\
\textbf{Decision-weight importance sampling (humans)} & \textbf{836} & \textbf{836} & \textbf{836} & \textbf{2,127} & \textbf{2,127} & \textbf{2,127} & \textbf{2,543} & \textbf{2,543} & \textbf{1,783}\\
Stratified by pilot variance (humans) & 823 & 823 & 823 & 2,076 & 2,076 & 2,076 & 2,482 & 2,482 & 1,785\\
Active inf., calibrated-judge rule + CV & 798 & 822 & 1,171 & 1,986 & 2,419 & 2,736 & 2,520 & 2,594 & 2,010\\
Active inf., residual-estimated rule + CV & 753 & 786 & 868 & 1,811 & 1,970 & 2,145 & 2,073 & 2,273 & 1,807\\
Active inf., robust mixture + CV & 748 & 776 & 886 & 1,820 & 2,013 & 2,182 & 2,011 & 2,192 & 1,822\\
Decision-weight sampling + judge CV ($\lambda=1$) & 762 & 794 & 879 & 1,788 & 2,004 & 2,120 & 2,084 & 2,278 & 1,789\\
\textbf{Decision-weight sampling + judge CV ($\lambda$ fitted)} & \textbf{770} & \textbf{772} & \textbf{822} & \textbf{1,720} & \textbf{1,838} & \textbf{2,024} & \textbf{2,046} & \textbf{2,049} & \textbf{1,722}\\\midrule
Pilot labels (common to every arm) & 631 & 631 & 631 & 975 & 975 & 975 & 1,029 & 1,029 & 1,387\\\bottomrule
\end{tabular}
\end{table}

\begin{figure}[t]\centering
\includegraphics[width=\linewidth]{figures/F7_J50_v3.png}
\caption{Cost of every arm relative to humans-only decision-weight sampling ($\Jfifty/\Jfifty^{w}$, precision at cutoff, $\eps=0.02$). Bars below one are cheaper than the human-only reference; the judge arms' increments are the small gaps below one, and among the judge-assisted arms the calibrated-judge rule with the adversarial judge is the only one far above it (the two uniform human-only arms lie above it by construction).}\label{fig:F7}
\end{figure}

We read Table~\ref{tab:J50} in four parts (paired-bootstrap 95\% intervals over the 300 draws in brackets). \textbf{The large saving needs no judge.} Decision-weight sampling needs $40\%$ [37, 44], $56\%$ [52, 59], $54\%$ [42, 58] and $56\%$ [53, 58] fewer labels than uniform sampling on ANTIQUE, CAsT, DBpedia-Entity and DL 21--23. A third human-only arm, uniform sampling restricted to the documents some comparison depends on, splits this: skipping the zero-weight documents saves $21\%$ [16, 26], $41\%$ [36, 45], $39\%$ [22, 45] and $36\%$ [31, 39], and weighting the rest by $|w|$ a further $25$--$32\%$ of that arm's cost. Stratifying by a pilot-estimated variance adds nothing beyond $|w|$. \textbf{The judge's increment is real but modest and conditional.} The $\lambda$-fitted control variate with Qwen3-8B lowers $\Jfifty$ by $8\%$ [5, 11], $19\%$ [14, 24], $20\%$ [13, 27] and $3\%$ [1, 6] (DL 21--23, where its $\rho\approx0.46$); the reranker by $8$--$19\%$; the adversarial judge by $2\%$ [$-2$, 4] and $5\%$ [$-1$, 9], nothing distinguishable from zero. The pilot, a fixed cost common to every arm, is $41$--$78\%$ of the weighted arm's $\Jfifty$ and dilutes the increment; in post-pilot labels it is $16$--$35\%$. \textbf{The judge can cost labels, and the fix is cheap.} With the adversarial judge the coefficient-1 control variate needs $5\%$ [2, 9] more labels than weighted sampling on ANTIQUE; fitting $\lambda\in[0,1]$ on the pilot removes the loss and is never worse than $\lambda=1$ with the good judges beyond noise. \textbf{No judge-assisted arm dominates.} The residual-estimated and robust active-inference rules and the two weighted control variates lie within $3$--$11\%$ of one another in every column; only the calibrated-judge rule fails badly ($+29$ to $+40\%$ over weighted sampling with the adversarial judge), because it withholds labels from documents on which the judge is confidently wrong, the failure anticipated by \citet{li2025robust}.

\paragraph{A cost model for the table, computable from the pilot.} With the population estimand and every non-pilot query audited, the variance behind the bound is the mean within-query document-sampling variance divided by the number of queries, which for $\pi_d\propto\mathrm{base}_d$ scales as $v_{\text{arm}}/L$ with $L$ the labels drawn after the pilot; the certificate crosses $\eps$ when $z\sqrt{v_{\text{arm}}/L}$ falls below the slack. Two arms sharing the pilot $P$, the candidate and the slack therefore satisfy $(\Jfifty^{a}-P)/(\Jfifty^{b}-P)\approx v_a/v_b$, and the total saving of $a$ over $b$ is $(1-v_a/v_b)(1-P/\Jfifty^{b})$: a variance reduction multiplied by the post-pilot share of the cost. Reading $v$ from the per-query variances at the smallest budgets, the model reproduces the 193 (collection, judge, $\eps$, arm) cells of Tables~\ref{tab:J50} and \ref{tab:J50b} with correlation $0.95$ between the logarithms of predicted and realised post-pilot cost ratios (median multiplicative error $1.10$) and the judge arms' total saving to a mean absolute error of $0.06$ (Figure~\ref{fig:F8}, Appendix~\ref{app:costmodel}); the same model fails on the menu-allocation designs (correlation $0.53$), which is why an oracle design buys little (Section~\ref{sec:negative}). Both factors are computable before any label beyond the pilot is bought: the within-query variance of every design on the pilot's fully labelled documents, and the post-pilot share from the pilot's slack. We recorded these pilot quantities in every draw and scored them against the realised costs of the same runs. Over the 18 judge cells at $\eps=0.02$ the pilot-predicted variance ratio correlates $0.95$ with the realised post-pilot cost ratio, and the fully pilot-based prediction of the total saving has mean absolute error $0.03$. A decision taken by the majority of pilots (adopt the judge if its predicted saving exceeds 5\%) is right in 15 of 18 cells, never recommends the adversarial judge, and every disagreement but one lies within two points of the threshold; the pilot ranks the $\lambda$-fitted above the coefficient-1 control variate in 9 of 9 cells. The post-pilot share is the weaker factor. The 20-query slack estimate predicts the reference arm's post-pilot labels only within a factor of two, so the rule answers whether a judge pays and which design is cheaper, but says less about how much.

\begin{figure}[t]\centering
\includegraphics[width=\linewidth]{figures/F8_cost_model.png}
\caption{The variance-dilution cost model on the fixed-budget runs. Left: post-pilot cost ratio of every arm to decision-weight sampling against the ratio of their within-query variances at the smallest budgets, 193 cells, $\eps\in\{0.01,0.02\}$; the dashed line is equality. Right: predicted total saving of the judge arms, $(1-v_a/v_w)(1-P/\Jfifty^{w})$, against the realised saving.}\label{fig:F8}
\end{figure}

\paragraph{A pre-registered test of the rule with a new judge.} We then tested the rule on the two fully judged pools that no document-level experiment had used, TREC-COVID (50 queries) and Touch\'e 2020 (49), at pilots of 20 and 10 queries, following the protocol of Appendix~\ref{app:heldout}: the pilot-only records and the predictions derived from them are committed before the audit, the audit is run with the same draws, and its pilot columns are checked to equal the committed ones. The test has two stages. With the reranker judge, only the label predictor above was committed in advance. Its post-pilot labels came out three to four times too high on these small pools ($565$ against $203$ realised on TREC-COVID, $244$ against $58$ on Touch\'e), because with a pilot of 20 of 50 queries the known pilot part of the population mean loosens the precision the rest must reach; we corrected the predictor after seeing those results (v2--v4 of Appendix~\ref{app:heldout}). With Qwen3-8B judgments generated afterwards for the same two pools, all four predictors were committed before the audit; the collections were no longer unseen, but the judge and every prediction were.

The cost ratios came out as the pilot had predicted: over all arms, predicted and realised post-pilot cost ratios correlate $0.96$--$0.97$ (reranker) and $0.94$--$0.95$ (Qwen3-8B). The corrected label predictor puts the post-pilot labels at $203$ and $47$. Its decisions, taken by the majority of the 300 pilots, were right in all eight pre-registered Qwen3-8B cells at pilot 20 (the original predictor in four): on TREC-COVID the judge saves $6$--$8\%$ and was recommended, on Touch\'e it saves $3$--$4\%$ and was not, its cost ratio being $0.38$ but the pilot $93\%$ of the cost. We designated this corrected predictor as the primary one after seeing these results; Appendix~\ref{app:heldout} reports all four. A majority over pilots is not what an auditor has, so we also score the executable rule: in each draw, use the $\lambda$-fitted control variate if that draw's own pilot predicts a saving above 5\%, otherwise sample without the judge. Individual pilots agree with the cell-level outcome in $52$--$99\%$ of draws on the held-out cells ($45$--$99\%$ on the development cells). The rule's cost lies between the two arms it chooses from (Table~\ref{tab:heldout}): on the development collections it keeps $75$--$100\%$ of the saving of always using the judge where the judge pays and never adopts the adversarial judge; on the held-out cells it gives up part of the saving at pilot 20 and matches always-using-the-judge at pilot 10, each draw charged its own pilot, with no wrong certificate. At pilot 10 every predictor recommends correctly in all eight cells while over-estimating its size (Table~\ref{tab:heldout}).

Two limits are visible already here: the absolute label count is predicted only to within a factor of about two, and single pilots disagree with one another near the threshold.

\begin{table}[t]\centering\footnotesize\setlength{\tabcolsep}{4pt}
\caption{Predictions fixed before the audit against realised outcomes, $\lambda$-fitted Qwen3-8B control variate over humans-only decision-weight sampling. ``Predicted'' is the locked v3 prediction of the total saving (median over pilots); ``always CV'' the realised saving of always using the judge; ``rule'' the realised saving of the single-pilot rule (adopt iff the draw's own pilot predicts more than 5\%); ``adopt'' the share of pilots that adopt. TREC-COVID and Touch\'e: predictions locked after their reranker audit; NeuCLIRBench: everything locked before any data were touched.}\label{tab:heldout}
\begin{tabular}{llrrrr rrrr}\toprule
 & & \multicolumn{4}{c}{$\eps=0.02$} & \multicolumn{4}{c}{$\eps=0.01$}\\
Collection & pilot & predicted & always CV & rule & adopt & predicted & always CV & rule & adopt\\\midrule
TREC-COVID & 20 & 6.3\% & 6.4\% & 5.7\% & 75\% & 7.9\% & 6.7\% & 6.3\% & 89\%\\
Touch\'e 2020 & 20 & 3.6\% & 2.7\% & 0.8\% & 19\% & 4.7\% & 3.5\% & 2.6\% & 48\%\\
TREC-COVID & 10 & 15.4\% & 14.8\% & 14.8\% & 91\% & 18.5\% & 15.7\% & 15.8\% & 91\%\\
Touch\'e 2020 & 10 & 12.0\% & 7.4\% & 7.5\% & 99\% & 16.3\% & 8.8\% & 8.8\% & 100\%\\\midrule
NeuCLIRBench & 20 & 4.5\% & 2.1\% & 1.6\% & 47\% & 5.7\% & 5.3\% & 2.3\% & 55\%\\
NeuCLIRBench & 10 & 7.0\% & 6.0\% & 4.6\% & 57\% & 8.2\% & 3.8\% & 5.0\% & 60\%\\\bottomrule
\end{tabular}
\end{table}

\paragraph{An independent test with the rule fixed in advance (Table~\ref{tab:heldout}).} Finally we fixed every element of the procedure, the corrected predictor, the threshold, the pilot size, the reference arm, the menu, the judges and the pool rule, in a lock committed before any data of the target had been touched, and applied it to NeuCLIRBench \citep{lawrie2025neuclirbench} in its monolingual English setting (105 topics with relevance judgments; documents are English machine translations of the Persian, Russian and Chinese news collections; pool of 71 judged documents per topic under the paper's pool rule; relevant $=$ gain $\ge1$), which had played no part in development. Qwen3-8B judgments were generated for the pool, the pilot-only predictions of all three judges were committed (Appendix~\ref{app:heldout}), and only then was the audit run. Judge accuracy against the gain-based labels is low here ($0.44$ for Qwen3-8B, $0.59$ for the reranker). The locked predictions were modest: the $\lambda$-fitted control variate with Qwen3-8B was predicted to save $4.5$--$8\%$ of the total and to be worth adopting in three of its four cells; both other judges were predicted not to pay.

The relative costs came out as the pilot had predicted: over all arms the correlation between predicted and realised post-pilot cost ratios is $0.83$--$0.92$ for Qwen3-8B and $0.89$--$0.97$ for the other judges, and the decomposition of the human-only saving (skipping zero-weight documents $0.62$, weighting $0.52$, together $0.82$ of the post-pilot labels) matches its prediction ($0.66$, $0.47$, $0.82$). The reranker and the inverted judge were correctly predicted not to pay in all eight cells, and the coefficient-1 control variate was correctly predicted to lose with both.

For Qwen3-8B the predictor was optimistic about the judge itself: its cost ratio was predicted at $0.71$--$0.76$ and realised at $0.82$--$0.93$, so the realised total savings fell below the predicted ones in every cell (Table~\ref{tab:heldout}). The locked majority decision was right in three cells and wrong in one (pilot 10, $\eps=0.01$, predicted $8.2\%$, realised $3.8\%$), but all four cells lie within three points of the threshold and the pilots themselves split almost evenly at pilot 20 ($47\%$ and $52\%$ of pilots recommending the judge), so the aggregate verdicts carry little information about what one pilot would decide. The optimism is consistent with fitting $\lambda$ and evaluating the residual variance on the same 20 pilot queries, which the locked predictor does not cross-fit, but the slack estimate and the shape of the certification curve can contribute as well, and we do not settle the attribution on this collection.

The single-pilot rule (Table~\ref{tab:heldout}, ``rule'') saved less than always using the judge in three of the four Qwen3-8B cells and more in one, with no wrong certificate above $0.013$ in any cell. Of the two usefulness criteria fixed in the lock, ``never more than two points below the better of the two always-policies'' was met in 9 of the 12 (judge, pilot, $\eps$) cells---the misses are Qwen3-8B at pilot 20 and $\eps=0.01$ ($2.3\%$ against $5.3\%$), the inverted judge at pilot 20 and $\eps=0.01$ ($0.1\%$ against $2.7\%$) and the reranker at pilot 10 and $\eps=0.01$ ($0.6\%$ against $2.8\%$), the latter two being gains of a judge that the rule was designed to ignore---and ``adopts the inverted judge in at most 10\% of pilots'' was met in one of four ($7$--$15\%$; these adoptions cost nothing because $\lambda$ is fitted near zero). The requirement that both hold in every setting was therefore not met. The post-pilot labels of the reference arm were under-predicted by $20$--$40\%$ ($324$ and $440$ predicted against $401$ and $592$ realised at pilot 20). The decision near the threshold is where a 20-query pilot is least informative.

\section{Experimental protocol}\label{sec:protocol}
\paragraph{Pools and policies.} Policies are cutoff rules over a shared pool: an expected-F1 gate with a target-calibrated scalar, a global score threshold, a learned truncation, and (for the self-preference test) a reranker-score threshold. Pools are built by the legacy protocol of the base study (BM25 plus three sentence-transformer retrievers, top-30 each, union) and a modern variant with Qwen3-Embedding-0.6B in place of the weakest retriever; the relevance model is a calibrated logistic regression on four pool features trained on four BEIR development collections and never on a target. \emph{Fully judged} pools take the NIST/BEIR-judged documents of each query as the corpus: TREC DL 2019 (43 queries) and 2020 (54), TREC-COVID (50), Touch\'e 2020 (49), DBpedia-Entity (399), and TREC DL 2021--2023 on MS~MARCO~v2 (211). The BEIR development pools have sparse judgments and are used only for the $\rho$ analysis.
\paragraph{Judges.} Qwen3-Reranker-0.6B (probability of ``yes''), Qwen3-8B and Mistral-7B-Instruct-v0.3 with the UMBRELA 0--3 prompt \citep{upadhyay2024umbrela} read from the next-token distribution over the digit tokens, an MPNet cosine rank rule (label-free, different family), and the inverted reranker as an adversarial control. Batched inference with left padding requires explicit position ids; without them one collection's judgments collapsed, which we detected and re-ran (Appendix~\ref{app:repro}).
\paragraph{Repeats and randomness.} 500 repeats for sequential certificates and 300 draws for fixed-budget designs; every method has its own random stream, so adding a method never perturbs another's numbers; all seeds are recorded.

\section{Pre-registered external evaluations}\label{sec:prereg}
\paragraph{TREC DL 2021--2023 (query-level methods).} Before downloading the data we committed a lock fixing the target, pooling rule, training data, $\eps=0.01$, the look grid (capped at $T=90$ by the 211 queries), the three methods (humans only, always-PPI, rule-based adoption), the judge, 500 repeats, and three criteria: validity (Clopper--Pearson upper limit of the wrong-certificate rate $\le0.15$), efficiency (always-PPI $\ge1.2\times$ humans-only in cumulative $\ACT$ at the largest look; rule-based $\ge$ humans-only), and adversarial rejection (the rule adopts the inverted judge in $\le5\%$ of runs). Validity and adversarial rejection passed for every method and judge (Table~\ref{tab:prereg}); efficiency did not: always-PPI reached $318$ against $313$ ($1.02\times$), rule-based $308$. The mechanism map explains why: the judge's $\rho$ is $0.40$--$0.52$, the unlabeled pool after auditing $121$ queries, and humans alone certify 63\% by $T=90$. This is below the $\rho$ at which Figure~\ref{fig:F6} first shows a gain, with a smaller unlabeled pool than any cell in which it did.

\begin{table}[t]\centering\small
\caption{Pre-registered evaluation on TREC DL 2021--2023 (500 repeats, looks $\le90$, $\eps=0.01$). Cumulative $\ACT$ out of 500 and wrong certificates.}\label{tab:prereg}
\begin{tabular}{llrrr}\toprule
Judge & Method & $\ACT$ by $T{=}50$ / $70$ / $90$ & wrong / 500 & adoption\\\midrule
--- & humans only & 84 / 198 / 313 & 0 & ---\\
Qwen3-8B & always-PPI & 105 / 220 / 318 & 0 & ---\\
Qwen3-8B & rule-based & 84 / 203 / 308 & 0 & 26\%\\
Qwen3-Reranker & always-PPI & 52 / 136 / 256 & 0 & ---\\
Qwen3-Reranker & rule-based & 84 / 194 / 309 & 0 & 1.3\%\\
inverted & always-PPI & 70 / 182 / 304 & 0 & ---\\
inverted & rule-based & 84 / 198 / 313 & 0 & 0\%\\
earlier in-sample certificate & --- & 434 / 494 / 498 & 7 & ---\\\bottomrule
\end{tabular}
\end{table}

\paragraph{TREC CAsT 2019 (document-level methods).} A second lock fixed a target never used before (159 conversational turns restricted to the judged MS~MARCO~v1 passages, 48.7 pooled documents per turn), the Qwen3-8B judge, a budget grid of $\{30,45,60,90\}$ full-query equivalents, $\eps\in\{0.01,0.02\}$, 300 draws, and four criteria: (P1) boundary type-I error $\le\alpha$ for every arm, (P2) $\Jfifty(\text{weighted})\le0.6\,\Jfifty(\text{uniform})$, (P3) the judge control variate never worse than weighted sampling for three judges, (P4) faithful active inference within 10\% of the weighted control variate. The judge reached pair accuracy $0.79$ and $\rho=0.62$, the most favourable judge conditions in this paper. P1 and P3 passed (maximum type-I error $0.013$). P2 and P4 could not be evaluated as written because \emph{no} arm reached 50\% $\ACT$ within the locked grid (best $0.43$): conversational queries make the policy gaps small, and the grid had been sized on easier collections. By the lock's own rule this counts as not met. The lesson concerns pre-registration rather than the methods: criteria must include a quantity that is always defined. Within the grid the ordering of Table~\ref{tab:J50} replicated exactly, and a post-hoc extension to 150 query-equivalents gave $\Jfifty$ of $2{,}750$ judged pairs for decision-weight sampling against \emph{not reached} for uniform sampling and $2{,}632$ with the 8B control variate (per-comparison accounting, Appendix~\ref{app:legacy}). The extension also showed that with the inverted judge the coefficient-1 control variate needs $16\%$ more judgments than weighted sampling alone: unbiasedness does not protect efficiency against an adversarial judge, as PPI without a tuned $\lambda$ would not. We therefore fit $\lambda\in[0,1]$ on the pilot and use $\lambda\hat\jmath$ as the control variate, so that $\lambda\to0$ recovers Horvitz--Thompson sampling. This change was made after the lock and became the first item of a third one.

\paragraph{ANTIQUE (the $\lambda$-corrected control variate).} The third lock targeted ANTIQUE \citep{hashemi2020antique}, a non-factoid question-answering test set never used in development (200 questions, 31.6 pooled answers per question), with a wider grid $\{30,60,90,120,150\}$ and criteria defined on $\ACT$ at the largest budget so that they are always evaluable. The Qwen3-8B judge, prompted on a 0--3 scale that does not match ANTIQUE's 1--4 scale, reached a pair accuracy of only $0.47$, yet $\rho=0.59$. Validity passed (boundary type-I error $\le0.010$); the $\lambda$-corrected control variate was never worse than weighted sampling for any judge and never worse than the coefficient-1 variate with the inverted judge, and the residual-estimated active rule matched it. The efficiency criterion ($\ACT(\text{weighted})\ge1.5\times\ACT(\text{uniform})$ at the largest budget) failed for the opposite reason to CAsT: the grid was large enough that every arm saturated ($0.96$ vs $0.92$), whereas the co-registered $\Jfifty$ read $776$ against $1{,}555$ ($0.50$). The lesson is symmetric to CAsT's: an always-defined criterion must be evaluated where the human-only arm is in its informative range, which we now fix by rule.

\section{Negative and bounded results}\label{sec:negative}
Four results bound the contribution. (1)~A band-only hybrid predictor (human labels on the disagreement band, judge elsewhere) raises $\rho$ from $0.6$ to $0.9$ for set-F1 but gives no gain at equal budget, because the rectifier still needs fully labelled queries; for precision it does not raise $\rho$ at all, as Proposition~\ref{prop:C} implies. (2)~Optimising the allocation of labels across the comparisons of a four-policy menu pays little, although \emph{sharing} them pays a lot: against independent per-comparison sampling with the same budget, static shared sampling needs $35$--$59\%$ fewer unique labels; an exact minimax design $\min_\pi\max_j \mathrm{Var}_j(\pi)/s_j^2$ solved with the true gaps and residual variances improves its objective $1.6$--$2.2\times$ over static sharing yet lowers $\Jfifty$ by only $0$--$14\%$ at $\eps=0.02$ and $1$--$17\%$ at $\eps=0.01$, and a plug-in version loses $13\%$ on DBpedia-Entity (Appendix~\ref{app:menu}). Two causes: in $8$--$16\%$ of draws the 20-query pilot picks a candidate whose true regret exceeds $\eps$, which caps every arm, and below the cap the oracle's advantage is confined to the budgets at which the conditional $\ACT$ is far from one. The cost model says the same thing: across the 29 menu cells the ratio of realised design objectives correlates only $0.53$ with the ratio of post-pilot costs, because $\max_j V_j/s_j^2$ averaged over draws is dominated by the draws with the smallest slack, which no allocation certifies, whereas $\Jfifty$ is set by the median draw. A larger pilot removes wrong candidates but its own price grows faster, so the total cost is minimised at 10--20 queries on every collection (Table~\ref{tab:pilot}). (3)~Two attempts to make always-PPI safe by adjusting $\lambda$ (finite-$N$ optimum, variance guard) reduced but did not remove its loss with a poor judge; the pilot adoption rule did. (4)~Judge-uncertainty sampling with a pilot-estimated residual model does not harm; only the calibrated-judge rule does (Table~\ref{tab:J50}; the history of this claim is in Appendix~\ref{app:repro}).

\section{Related work}\label{sec:related}
\paragraph{Prediction-powered and active inference.} PPI \citep{angelopoulos2023ppi} and PPI++ \citep{angelopoulos2023ppipp} correct AI predictions with human labels and tune $\lambda$ to bound the loss from poor predictions; active statistical inference \citep{zrnic2024active} chooses which points to label while keeping the guarantee; \citet{li2025robust} mix uncertainty-based sampling with a safe distribution; \citet{sfyraki2026revisiting} find constant-probability sampling competitive for sequential mean estimation; \citet{kilian2025anytime} give anytime-valid PPI. We use this machinery unchanged: bias removal by a control variate is PPI's, and the allocation of labels by residual variance is active inference's. What we add is the application of both to the simultaneous certification of a menu of retrieval policies (the functional the problem requires, and the union bound that keeps PPI's judge-independence under data-driven candidate choice), a cost accounting that charges every human label including the pilot, and a measurement of how far a pilot can foresee that cost.
\paragraph{AI judges for IR evaluation.} \citet{oosterhuis2024reliable} build PPI and conformal intervals for IR metrics from LLM judgments; \citet{chatzi2024ppr} rank LLMs from LLM comparisons with human anchoring; \citet{gligoric2024unconfident} use LLM confidence to allocate human annotation; \citet{balog2025rankers} document that LLM judges favour LLM-based rankers; UMBRELA \citep{upadhyay2024umbrela} is the source of our prompt; \citet{durmazkeser2026selection} select among many LLMs with few labels. These treat the judge as a fixed predictor of a metric or ranking. We add that for deployment decisions the judge's value is set by $\rho$ of paired differences, not accuracy, that this varies strongly by judge family and collection, and the measured self-preference of the identical-model judge.
\paragraph{Sequential testing, best-arm identification, IR evaluation cost.} Adaptive learn-then-test \citep{zecchin2024altt}, fixed-confidence BAI \citep{jamieson2014best,kaufmann2016complexity}, $\eps$-best-answer identification \citep{jourdan2022choosing} and multi-task PPI \citep{emmenegger2026multitask} give stopping rules and shared predictors; we use a plain union bound over looks and ordered pairs and handle the winner's curse of data-driven candidate choice explicitly. Topic-set size design \citep{sakai2016topic}, few-topic prediction \citep{guiver2009few}, active judgment sampling \citep{li2017active}, active comparison of two models \citep{sawade2012active} and active testing \citep{kossen2021active} trade labels for evaluation precision; we inherit Sawade et al.'s principle of labelling where the difference lives and carry it to the document level with a certification target and a cost accounting.

\section{Limitations and conclusion}
Our pools consist of judged documents only, so conclusions concern deployment decisions over judged candidate sets. Precision at cutoff is the primary utility for the document-level results; set-F1 is treated by a bias-corrected linearisation whose remainder we measure (Appendix~\ref{app:f1}). The certificate uses Bonferroni over looks rather than an anytime-valid sequence, and its $t$ bound is calibrated only asymptotically: the wrong-certificate rates we report are observed over 300--500 repeats per cell, not guarantees, and the exact bound we tried is affordable on two collections and close to a census on two others. The cost figures are fixed-budget certification costs, not sequential stopping costs; they charge every unique label once, including a 20-query pilot that is $41$--$78\%$ of the weighted arm's $\Jfifty$, and the population sizes ($159$--$399$ queries) bound how far any budget curve can be followed. The pre-audit rule has been tested under a lock with a new judge on two collections of 49--50 queries whose reranker results had already been examined, and, with the rule fixed in advance, on one collection of 105 queries never used in development, where every judge was weak and the pilot was $70$--$78\%$ of the cost; a locked test on a collection with a hard decision and a high-$\rho$ judge, where the judge's increment and the rule's value should be largest, has not been done. Three judge families were used, with one LLM judge per family. Within these limits the conclusion is simple. Once the estimand that certification requires is defined, existing tools (importance sampling by decision weight, a judge as control variate, a split certificate) deliver most of what is available. What a design is worth is its variance reduction times the post-pilot share of the cost; a 20-query pilot can tell the auditor the relative cost of designs before the audit, and the judge's own increment less reliably.

\bibliographystyle{tmlr}
\bibliography{refs}

\appendix
\section{Set-F1: linearised document-level auditing}\label{app:f1}
For set-F1, $D(y)=2\,\mathrm{tp}_a/(k_a+n_G)-2\,\mathrm{tp}_b/(k_b+n_G)$ is a ratio of label sums. We linearise at the judge labels: $w_d=\partial D/\partial y_d|_{y=\hat\jmath}=2\cdot\mathbf 1[d\in R_a]/(k_a+\hat n_G)-2\cdot\mathbf 1[d\in R_b]/(k_b+\hat n_G)-2\hat{\mathrm{tp}}_a/(k_a+\hat n_G)^2+2\hat{\mathrm{tp}}_b/(k_b+\hat n_G)^2$, so every pooled document has a non-zero weight through $n_G$; $\hat D_{\mathrm{lin}}=D(\hat\jmath)+\sum_{d\text{ sampled}}w_d(y_d-\hat\jmath_d)/\pi_d$ with $\pi_d\propto|w_d|$. The remainder $D-\hat D_{\mathrm{lin}}$ is measured on the fully labelled pilot and its mean is added as a bias correction with its standard error folded into the bound. A humans-only plug-in alternative estimates $\mathrm{tp}_a,\mathrm{tp}_b,n_G$ by Horvitz--Thompson and plugs them into F1.

Measured remainders: DBpedia-Entity $+0.002$ (mean absolute $0.011$), DBpedia with the reranker judge $-0.012$ ($0.018$), TREC DL 21--23 $-0.029$ ($0.030$), i.e.\ three times $\eps=0.01$ on the last collection, where the uncorrected linearised certificate is not trustworthy. At equal document budget (60 full-query equivalents, $\eps=0.01$, 300 draws) the $\ACT$ rates on DBpedia-Entity were: full-query audit $0.23$, plug-in HT (humans) $0.78$, query-level PPI $0.33$, linearised CV $0.55$, bias-corrected linearised CV $0.32$; on TREC DL 21--23 (30 full-query equivalents): $0.10$, $0.72$, $0.10$, $0.86$ (inflated by the remainder), $0.42$. The boundary stress test of Section~\ref{sec:where} gave type-I error $\le0.013$ for all set-F1 arms at realised $\eps\approx0.03$; it does not probe $\eps=0.01$ on TREC DL 21--23, where only the bias-corrected variant should be used.

\section{Proofs}\label{app:proofs}
\paragraph{Proposition~\ref{prop:A}.} The two sample means are independent; $\mathrm{Var}[\overline{(D-\lambda\hat D)}_n]=(\sigma^2-2\lambda\rho\sigma\hat\sigma+\lambda^2\hat\sigma^2)/n$ and $\mathrm{Var}[\lambda\overline{\hat D}_N]=\lambda^2\hat\sigma^2/N$. The sum is a quadratic in $\lambda$ with derivative $(-2\rho\sigma\hat\sigma+2\lambda\hat\sigma^2)/n+2\lambda\hat\sigma^2/N$, which vanishes at $\lambda^\star=\rho\sigma/(\hat\sigma(1+n/N))$; substituting gives $\sigma^2/n-\rho^2\sigma^2/(n(1+n/N))=\frac{\sigma^2}{n}\big(1-\frac{\rho^2}{1+n/N}\big)$, and $N\to\infty$ recovers $\lambda^\star=\rho\sigma/\hat\sigma$ and $\sigma^2(1-\rho^2)/n$.\qed
\paragraph{Proposition~\ref{prop:B}.} Given $\mathcal G_t$ (the training half and its decisions), the validation queries are an i.i.d.\ sample from $P$ independent of $\mathcal G_t$. For the human-only bound this is the usual $t$ bound. For the PPI bound, within each fold $\lambda$ is estimated on the other fold, hence independent of the fold it is applied to, and the unlabeled mean is taken outside $B_t$; conditionally on $\lambda$, for a fixed ordered pair $(j,m)$ with $D_{jm}=u_j-u_m$ and $\hat D_{jm}=\hat u_j-\hat u_m$, $\mathbb E[D_{jm}-\lambda\hat D_{jm}]=\mu_j-\mu_m-\lambda\mathbb E[\hat D_{jm}]$ and $\mathbb E[\lambda\overline{\hat D}_{jm,N}]=\lambda\mathbb E[\hat D_{jm}]$, so the estimator of every fixed contrast is conditionally unbiased for any $f$; conditionally on the fold's $\lambda$ its variance is that of Proposition~\ref{prop:A} at that $\lambda$. The variance of the combined cross-fitted estimator, which also depends on the variation of the estimated $\lambda$ and on the dependence between folds, is not bounded by the proposition; the standard error the implementation uses (residual sample variance plus the variance of the unlabeled mean) and the certificate built on it were evaluated only empirically, in the wrong-certificate rates reported. (Unbiasedness is asserted for fixed contrasts only; the estimator of the contrast selected on $B_t$ is in general not unbiased, and nothing below requires it to be.) Let $\mathcal E$ be the event that every one of the $L|\mathcal M|(|\mathcal M|-1)$ bounds $U_{t,j,m}$ (all looks, all ordered pairs of the menu fixed before the audit) covers its target; by the union bound $\Pr(\mathcal E^c)\le L|\mathcal M|(|\mathcal M|-1)\alpha'=\alpha$ whenever each bound has level $1-\alpha'$ (exactly for the empirical-Bernstein and betting bounds on bounded differences; asymptotically for the $t$ bound). On $\mathcal E$, for any candidate $\hat m_t\in\mathcal M$, however it was chosen, the pair $(j^\star,\hat m_t)$ with $j^\star=\arg\max_j\mu_j$ belongs to the covered family, so $\mu_{j^\star}-\mu_{\hat m_t}\le U_{t,j^\star,\hat m_t}$; if in addition $\ACT_t$ holds then $U_{t,j^\star,\hat m_t}\le\eps$, so the candidate is not $\eps$-wrong. Hence a wrong certificate at any look implies $\mathcal E^c$. Choosing $\hat m_t$ on $A_t$ (an option of the implementation; the reported runs choose it on $B_t$) is therefore not needed for validity; it matters only if the family of contrasts could be enlarged after inspecting $B_t$, which the fixed menu rules out.\qed
\paragraph{Proposition~\ref{prop:C}.} Documents in $R_a\cap R_b$ contribute $y_d/Z_a-y_d/Z_b$, documents in $R_b\setminus R_a$ contribute $-y_d/Z_b$, and documents outside $R_a\cup R_b$ contribute nothing.\qed

\section{Pre-registered cost table under the original accounting}\label{app:legacy}
Table~\ref{tab:J50legacy} is the document-level cost table as it stood in our pre-registered runs. Its accounting differs from Table~\ref{tab:J50} in three ways: documents were averaged per comparison rather than counted once per menu, labels shared by the three comparisons were not de-duplicated (and the per-pair arm drew its own sample for each comparison, which a label ledger later showed cost it more than twice the unique labels of the shared arms), and the pilot was charged only up to the largest policy cutoff rather than for its full pools. The ordering of methods is the same under both accountings; the size of the judge increment is not, because the pilot is a fixed cost common to every arm. We keep the table as the historical record and draw no cost conclusion from it.

\begin{table}[h]\centering\footnotesize
\caption{Original per-comparison accounting (precision at cutoff, $\eps=0.02$, $\alpha=0.10$, 300 draws; observed wrong-certificate rate $\le0.007$). Judge: Qwen3-8B unless stated. Superseded by Table~\ref{tab:J50}.}\label{tab:J50legacy}
\begin{tabular}{lrrr}\toprule
Method & DBpedia-Entity & DBpedia (reranker judge) & TREC DL 21--23\\\midrule
Uniform document sampling (humans) & n.r.\ ($>2{,}570$) & n.r. & 3{,}207\\
Decision-weight importance sampling (humans) & 1{,}990 & 1{,}990 & 1{,}273\\
Stratified by pilot variance (humans) & 2{,}144 & 2{,}144 & 1{,}388\\
Active inference, calibrated-judge rule + CV & 2{,}020 & 2{,}054 & 1{,}562\\
Active inference, residual-estimated rule + CV & 1{,}610 & 1{,}715 & 1{,}340\\
Active inference, robust mixture + CV & 1{,}556 & 1{,}628 & 1{,}311\\
Decision-weight sampling + judge control variate & 1{,}569 & 1{,}705 & 1{,}256\\\bottomrule
\end{tabular}
\end{table}


\begin{table}[h]\centering\scriptsize\setlength{\tabcolsep}{3pt}
\caption{Table~\ref{tab:J50} at $\eps=0.01$ (same accounting, estimand and grids). ``n.r.''\ = 50\% $\ACT$ not reached at the largest budget run ($\ACT$ 0.42 on CAsT and 0.38 on DBpedia-Entity). Observed wrong-certificate rate $\le0.013$.}\label{tab:J50b}
\begin{tabular}{l rrr rrr rr r}\toprule
 & \multicolumn{3}{c}{ANTIQUE} & \multicolumn{3}{c}{CAsT 2019} & \multicolumn{2}{c}{DBpedia-Entity} & DL 21--23\\
Method & Q & R & I & Q & R & I & Q & R & Q\\\midrule
Uniform document sampling (humans) & 1,576 & 1,576 & 1,576 & n.r.\ ($>$5,683) & n.r. & n.r. & n.r.\ ($>$6,879) & n.r. & 4,579\\
Uniform over decision-relevant documents (humans) & 1,220 & 1,220 & 1,220 & 3,485 & 3,485 & 3,485 & 5,839 & 5,839 & 3,189\\
Decision-weight importance sampling (humans) & 851 & 851 & 851 & 2,775 & 2,775 & 2,775 & 3,581 & 3,581 & 1,958\\
Stratified by pilot variance (humans) & 877 & 877 & 877 & 2,737 & 2,737 & 2,737 & 3,870 & 3,870 & 1,979\\
Active inf., calibrated-judge rule + CV & 807 & 884 & 1,273 & 2,459 & 2,906 & 3,384 & 3,979 & 4,362 & 2,298\\
Active inf., residual-estimated rule + CV & 781 & 814 & 925 & 2,340 & 2,500 & 2,726 & 2,859 & 3,345 & 1,896\\
Active inf., robust mixture + CV & 784 & 812 & 943 & 2,324 & 2,478 & 2,682 & 2,945 & 3,151 & 1,869\\
Decision-weight sampling + judge CV ($\lambda=1$) & 794 & 817 & 934 & 2,320 & 2,657 & 2,746 & 2,959 & 3,456 & 1,917\\
Decision-weight sampling + judge CV ($\lambda$ fitted) & 801 & 803 & 870 & 2,190 & 2,421 & 2,633 & 3,237 & 3,347 & 1,883\\\midrule
Pilot labels (common to every arm) & 632 & 632 & 632 & 975 & 975 & 975 & 1,028 & 1,028 & 1,387\\\bottomrule
\end{tabular}
\end{table}

\section{Allocation across the comparisons of a menu}\label{app:menu}
For a menu with candidate $\hat m$ and $J=|\mathcal M|-1$ competitors, a document-level design is one inclusion probability $\pi_d$ per document, shared by every comparison, under a budget $\sum_d c_d\pi_d\le B$. The value of sharing itself is measured by a per-comparison arm that draws an independent sample of $B/J$ documents for each comparison (duplicates across comparisons charged once by the ledger): its $\Jfifty$ at $\eps=0.02$ is $1{,}784$ / n.r.\ ($\ACT\le0.06$ at $4{,}486$) / $5{,}308$ / $3{,}234$ unique labels on ANTIQUE / CAsT / DBpedia-Entity / DL 21--23 against $1{,}015$ / $3{,}181$ / $2{,}186$ / $2{,}102$ for static shared sampling, i.e.\ sharing saves $35$--$59\%$ where both are measurable. The remainder of this appendix asks how much \emph{optimising} the shared allocation adds on top. Comparison $j$ has estimator variance $V_j(\pi)=\sum_d a_{jd}^2 v_d(1-\pi_d)/\pi_d$ ($a_{jd}$ its decision weight, $v_d$ the residual variance of the control variate) and slack $s_j=\eps-(\mu_j-\mu_{\hat m})$; the certificate fires when every $V_j$ is small relative to $s_j^2$, so the natural design is $\min_\pi\max_j V_j(\pi)/s_j^2$. Its dual over the simplex has a closed-form water-filling inner solution, $\pi_d=\min\{1,\sqrt{g_d/(\mu c_d)}\}$ with $g_d=v_d\sum_j\theta_j a_{jd}^2/s_j^2$, and we solve the outer problem by exponentiated gradient. Two versions were run: an \emph{oracle} that knows the true slacks and the population residual-variance function (but not the labels), and a \emph{plug-in} that estimates both from the pilot and the labels accumulated so far. Both were compared with static shared sampling $\pi_d\propto\sum_j|a_{jd}|$ (the arm of Table~\ref{tab:J50}), its variance-aware variant $\pi_d\propto\sum_j|a_{jd}|\sqrt{\hat v_d}$, and an adaptive re-weighting of comparisons by their estimated slack, all with the same ledger, pilot, and certificate, on the four target collections at 300 draws.

\begin{table}[h]\centering\scriptsize\setlength{\tabcolsep}{3pt}
\caption{Menu allocation, $\eps=0.02$ and $0.01$, Qwen3-8B judge, population estimand, 300 draws. ``Design ratio'' is the realised $\max_jV_j/s_j^2$ of static shared sampling divided by the oracle's (structural head-room, independent of the bound); ``saving'' is $1-\Jfifty(\text{arm})/\Jfifty(\text{static})$ in unique labels (Monte-Carlo error 6--9\%); ``infeasible'' is the fraction of draws in which the pilot's candidate has true regret $>\eps$; ``$\ACT\mid$feasible'' is the certification rate among the other draws at the largest budget run; ``non-overlap'' is the share of the binding comparison's decision weight on documents that no other comparison weights. Observed wrong-certificate rate $\le0.003$.}\label{tab:menu}
\begin{tabular}{llrrrrrr}\toprule
Collection & $\eps$ & design ratio & saving (oracle) & saving (plug-in) & infeasible & $\ACT\mid$feas.\ (static / oracle) & non-overlap\\\midrule
ANTIQUE & 0.02 & 1.6 & 2\% & 2\% & 0.15 & 0.98 / 1.00 & 0.59\\
CAsT 2019 & 0.02 & 2.0 & 14\% & 12\% & 0.08 & 1.00 / 1.00 & 0.60\\
DBpedia-Entity & 0.02 & 1.6 & 0\% & $-13\%$ & 0.08 & 0.85 / 0.89 & 0.56\\
TREC DL 21--23 & 0.02 & 2.2 & 7\% & 2\% & 0.13 & 0.98 / 1.00 & 0.83\\\midrule
ANTIQUE & 0.01 & 1.6 & 1\% & $-1\%$ & 0.14 & 1.00 / 1.00 & 0.58\\
CAsT 2019 & 0.01 & 2.2 & 17\% & 13\% & 0.15 & 0.97 / 1.00 & 0.60\\
DBpedia-Entity & 0.01 & 1.6 & 1\% & $-12\%$ & 0.16 & 0.83 / 0.87 & 0.55\\
TREC DL 21--23 & 0.01 & 2.1 & 6\% & 3\% & 0.13 & 1.00 / 1.00 & 0.83\\\bottomrule
\end{tabular}
\end{table}

Three facts explain the gap between the design ratio and the saving, and they do not weigh equally on every collection. (i)~\emph{The pilot's candidate choice caps every arm.} With a 20-query pilot, $8$--$16\%$ of draws select a candidate whose true regret exceeds $\eps$; no allocation can certify these correctly. On ANTIQUE, CAsT and DL 21--23 the conditional $\ACT$ among feasible draws reaches $0.97$--$1.00$ for both arms at the largest budget, so the observed plateau is this cap and nothing else. On DBpedia-Entity the conditional $\ACT$ is $0.83$--$0.89$ at 150 query-equivalents and still rising; the pilot-size and fixed-candidate runs below show that this residual is the small slack of feasible candidates relative to the bound's width, not the pilot. (ii)~\emph{The oracle's advantage lives only in the transition.} It is largest where the conditional $\ACT$ is far from its cap and the budget is binding: CAsT at $\eps=0.01$ ($+0.31$ and $+0.32$ in conditional $\ACT$ at 90 and 120 query-equivalents, which is the $17\%$ saving), DL 21--23 at 5--10 query-equivalents ($+0.10$--$0.19$); elsewhere it is below $0.05$. Two facts about CAsT make it the collection where the design matters: its policies are close (median slack $\approx\eps$, $7$--$15\%$ of feasible draws with slack below $\eps/2$) and its comparisons' decision weights overlap little (non-overlap $0.60$), so concentrating labels on the binding comparison's documents is both needed and possible. (iii)~\emph{The pilot dilutes the saving.} The pilot is $30$--$50\%$ of the labels at $\Jfifty$ for these arms and is identical across them, so a proportional saving on the sampled labels appears as a smaller saving on the total. Under the earlier version of these runs, which bounded the mean over the non-pilot queries rather than the population mean, CAsT did not certify at all at these budgets; with the estimand aligned it does, and the ranking of the arms is unchanged.

\paragraph{Pilot size and the candidate-selection effect.} Because the pilot both selects the candidate and is paid for, we varied it ($10$, $20$, $40$, $80$ fully labelled queries; the candidate chosen on the pilot is shared by every arm within a draw) and, separately, removed the selection effect by forcing every draw to certify the true best policy (regret $0$, slack $\eps$) with a 20-query pilot. Table~\ref{tab:pilot} reports the static shared arm; the oracle and plug-in arms move with it within $0$--$14\%$ at every pilot size, so a better candidate does not unmask a hidden allocation gain. Three facts follow. (a)~Larger pilots remove wrong candidates: the share of draws with a certifiable candidate rises from $0.71$--$0.87$ (pilot 10) to $0.99$--$1.00$ (80) at $\eps=0.02$. (b)~Given a certifiable candidate, the shared arms certify (conditional $\ACT\ge0.95$ on ANTIQUE, CAsT and DL 21--23 at every pilot size); DBpedia-Entity is the exception ($0.85$--$0.90$), and with the true best candidate it reaches $0.99$ at $\eps=0.02$ but only $0.90$ at $\eps=0.01$, so its residual is feasible candidates with small slack, not the pilot. (c)~Under full accounting a larger pilot never pays: the labels drawn after the pilot fall with its size, but its own labels grow faster, so the total is lowest at pilot 10 or 20 everywhere, $1.1$--$1.7\times$ higher at 40 and $1.3$--$3.1\times$ at 80. Removing the selection effect (true best candidate, pilot 20) lowers the total by $2\%$ on ANTIQUE and DL 21--23 and by $18$--$21\%$ on DBpedia-Entity and CAsT at $\eps=0.02$, where the policies are close. We pre-registered three predictions for $\eps=0.03$, a tolerance not used in any development run, before running it (repository, review response \S14): monotone cost in pilot size above 20 (held on 12/12 arm--collection cells); conditional $\ACT\ge0.95$ on the three non-DBpedia collections (held); an oracle saving $\le15\%$ (held, $0$--$4\%$). The same prediction also claimed pilot 10 would be within $\pm15\%$ of pilot 20, and this part \emph{failed}: at $\eps=0.03$ pilot 10 is $3$--$13\%$ cheaper than pilot 20 on CAsT and $18$--$30\%$ cheaper on the other three collections, i.e.\ the looser the tolerance, the smaller the pilot that minimises total cost, because a wrong candidate becomes rarer while the pilot's price does not change.

\begin{table}[h]\centering\scriptsize\setlength{\tabcolsep}{3pt}
\caption{Pilot size, static shared arm, population estimand, $\eps=0.02$, 300 draws per budget. ``Feasible'' is the share of draws whose pilot-chosen candidate has true regret $\le\eps$; ``$\ACT\mid$feas.'' the certification rate among those draws at the largest budget; ``after pilot'' and ``total'' are the unique human-labelled documents to reach 50\% $\ACT$, excluding and including the pilot. Wrong-certificate rate $\le0.003$ in every cell.}\label{tab:pilot}
\begin{tabular}{lrrrrr}\toprule
Collection & pilot & feasible & $\ACT\mid$feas. & after pilot & total\\\midrule
ANTIQUE & 10 & 0.73 & 0.96 & 574 & 890\\
 & 20 & 0.85 & 0.98 & 381 & 1{,}015\\
 & 40 & 0.93 & 1.00 & 320 & 1{,}585\\
 & 80 & 0.99 & 1.00 & 240 & 2{,}769\\
\midrule
CAsT 2019 & 10 & 0.83 & 0.98 & 2{,}561 & 3{,}050\\
 & 20 & 0.92 & 1.00 & 2{,}205 & 3{,}181\\
 & 40 & 0.97 & 1.00 & 1{,}665 & 3{,}618\\
 & 80 & 1.00 & 1.00 & 879 & 4{,}779\\
\midrule
DBpedia-Entity & 10 & 0.87 & 0.88 & 1{,}095 & 1{,}605\\
 & 20 & 0.92 & 0.85 & 1{,}155 & 2{,}186\\
 & 40 & 0.96 & 0.89 & 744 & 2{,}797\\
 & 80 & 1.00 & 0.90 & 636 & 4{,}749\\
\midrule
TREC DL 21--23 & 10 & 0.71 & 0.95 & 1{,}056 & 1{,}753\\
 & 20 & 0.87 & 0.98 & 717 & 2{,}102\\
 & 40 & 0.98 & 1.00 & 509 & 3{,}283\\
 & 80 & 1.00 & 1.00 & 325 & 5{,}859\\\bottomrule
\end{tabular}
\end{table}

We conclude that, at these population sizes, the certification cost is governed first by the pilot---its price more than its candidate choice on ANTIQUE and DL 21--23, both on CAsT and DBpedia-Entity---and by the decision margin, and only then by the allocation of labels across comparisons, with CAsT as the one collection where the allocation is worth its cost. The practical reading is not ``choose a larger pilot'': it is to diagnose the candidate's failure rate and slack before optimising where the remaining labels go.

\section{Exact and finite-population certificates}\label{app:exact}
The estimand is fixed to the mean over the $N$ queries of the evaluation set; the training half $A$ of the audit is fully labelled and therefore known exactly, and the validation half $B$ is a uniformly random sample without replacement from the remaining $N-|A|$ queries, so that $\mu_N=(\sum_{q\in A}D_q+(N-|A|)\,\mu_{\text{rest}})/N$ and only $\mu_{\text{rest}}$ needs a bound. Three bounds on $\mu_{\text{rest}}$ were compared for the human-only certificate: the pre-registered one-sided $t$ bound on $B$ alone (targeting the superpopulation mean, no use of $A$); the same $t$ bound with the finite-population correction $\sqrt{1-|B|/(N-|A|)}$ inside the decomposition above; and the exact without-replacement betting bound of \citet{waudby2020wor}, in which the hypothesis ``$\mu_{\text{rest}}=m$'' is tested by a nonnegative supermartingale whose predictable conditional mean after $i$ draws is $(N'm-S_i)/(N'-i)$, and hypotheses that become impossible given the observed prefix are rejected outright. The judge-assisted version bounds the rectifier $D_q-\lambda\hat D_q$ (range $[-1-\lambda,1+\lambda]$, $\lambda\in[0,1]$ fitted on $A$) and adds the exactly known population mean of $\lambda\hat D_q$. The training half was fixed at 20 queries so that the validation half could grow to $N-20$; every arm was evaluated at every look until it certified or the last look was reached (the last look audits $N-1$ queries). Level $\alpha'=\alpha/(L|\mathcal M|(|\mathcal M|-1))$ with $L=14$ looks, $\eps=0.01$, 300 repeats; wrong certificates were scored against the population mean. $T$ counts all audited queries, the 20-query training half included, so $T/N$ is the share of the population labelled.

\begin{table}[h]\centering\scriptsize\setlength{\tabcolsep}{2pt}
\caption{Split certificates under three bounds, population estimand, $\eps=0.01$, 300 repeats. $t$ columns: certification rate and mean stopping look $\bar T$ (queries audited, pilot included). Betting columns: cumulative certification rate by the time $75\%$ / $90\%$ / $100\%$ of the population has been audited, and $\bar T$ for the human-only arm. The last column is the largest look run. Observed wrong-certificate rate $\le0.007$ ($t$) and $0$ (betting) in every cell.}\label{tab:exact}
\begin{tabular}{lrrrrrr}\toprule
Collection ($N$) & $t$, pre-reg. & $t$+FPC+known $A$ & $t$+FPC+judge & bet.\ WoR (75/90/100\%) & bet.\ WoR+judge & max $T/N$\\\midrule
ANTIQUE (200) & $0.65$ ($\bar T{=}148$) & $1.00$ ($\bar T{=}110$) & $1.00$ ($\bar T{=}103$) & $0.02$ / $0.03$ / $0.88$ ($\bar T{=}198$) & $0.00$ / $0.01$ / $0.71$ & 100\%\\
CAsT 2019 (159) & $0.77$ ($\bar T{=}121$) & $1.00$ ($\bar T{=}88$) & $1.00$ ($\bar T{=}79$) & $0.00$ / $0.05$ / $0.91$ ($\bar T{=}152$) & $0.00$ / $0.01$ / $0.73$ & 99\%\\
DBpedia-Entity (399) & $0.94$ ($\bar T{=}142$) & $1.00$ ($\bar T{=}113$) & $1.00$ ($\bar T{=}96$) & $0.90$ / $0.95$ / $1.00$ ($\bar T{=}218$) & $0.56$ / $0.94$ / $1.00$ & 100\%\\
TREC DL 21--23 (211) & $0.99$ ($\bar T{=}92$) & $1.00$ ($\bar T{=}71$) & $1.00$ ($\bar T{=}68$) & $0.73$ / $0.73$ / $1.00$ ($\bar T{=}161$) & $0.47$ / $0.47$ / $1.00$ & 100\%\\\bottomrule
\end{tabular}
\end{table}

A synthetic check isolates the cause. For $N=399$, paired differences with mean $-0.02$ and standard deviation $0.05$ in $[-1,1]$, and $\alpha'\approx10^{-3}$, the betting bound at $n=100$ is $+0.06$ while the $t$ bound is $-0.004$; the betting bound crosses $\eps=0.01$ only at $n\approx250$ ($63\%$ of the population), and this is insensitive to the resolution of the candidate-mean grid. With range $R=2$ the second-order term $R\log(1/\alpha')/n$ of any empirical-Bernstein-type bound is $0.14$ at $n=100$, fourteen times $\eps$; halving the range would bring the crossing to $n\approx100$. The runs of Table~\ref{tab:exact} use the set-F1 utility of the query-level block and the loose range $[-1,1]$ for every paired difference. A narrower range is in fact known before any label is seen: for precision at cutoff with nested cutoffs $k_a\le k_b$ the paired difference of a query lies in $\pm(1-k_a/k_b)$ (e.g.\ $[-0.5,0.5]$ for $k_a=10$, $k_b=20$), whereas for set-F1 the range depends on $n_G$ and is not known label-free. Neither was used, so Table~\ref{tab:exact} prices \emph{this} bound with the loose range, and a range-aware betting bound is the obvious next variant. Dropping the union bound over looks, which a betting sequence does not need, lowers $\log(1/\alpha')$ from $7.1$ to $4.8$ and the width by about $15\%$. On ANTIQUE and CAsT the betting certificate fires almost only at the last look ($N-1$ queries audited), where the unobserved remainder is a single query; whether the difference from DBpedia-Entity and DL 21--23 is the slack of the binding comparison, the population size or the residual variance is not settled by these runs. The finite-population $t$ certificate is a strict improvement on the pre-registered one at no cost in observed error and is recommended whenever the guarantee is meant for the fixed evaluation set, never as a guarantee about a future query distribution.

\section{The variance-dilution cost model}\label{app:costmodel}
Section~\ref{sec:table2} states the model; here is its test, preceded by the estimand it rests on. The bound of the fixed-budget certificate targets the mean paired difference over all $N$ queries: the pilot's contribution is known exactly, and the remaining queries are a without-replacement sample whose per-query document-level estimates carry the two-stage variance $(1-f)\,s_b^2/n+f\,\overline{\hat v}/n$, with $f$ the fraction of the remaining queries audited and $\hat v$ the Horvitz--Thompson variance of each query's estimate; when every query is audited only document-sampling noise remains. An earlier version bounded the mean over the non-pilot queries as an i.i.d.\ sample while scoring against the population mean; on the same draws, decision-weight sampling on CAsT needs $3{,}907$ labels under the i.i.d.\ bound, $2{,}600$ once the pilot is known exactly, and $2{,}127$ with the two-stage variance, against $881\to851\to836$ on ANTIQUE, the difference being the pilot's share of the population ($12.6\%$ on CAsT). For every (collection, judge, $\eps$) cell of the fixed-budget runs and every arm $a$, we compare the realised post-pilot cost ratio $(\Jfifty^{a}-P)/(\Jfifty^{w}-P)$, with $w$ decision-weight sampling, to the ratio of the arms' per-query estimator variances (the \texttt{var\_est} column of the result files) averaged over the three smallest budgets, where $b=4$ documents per query and the within-query sampling variance dominates the between-query variance that every arm shares. The predicted total saving of a judge arm is $(1-v_a/v_w)(1-P/\Jfifty^{w})$. Figure~\ref{fig:F8} shows both comparisons; the numbers are produced by \texttt{95\_cost\_model.py} from the committed result files. Over 193 cells the log cost ratio and the log variance ratio correlate at $0.95$; the median multiplicative error is $1.10$ and the 90th percentile $1.37$. By arm, the median absolute log error is $0.05$--$0.11$ except for the two arms that sample by $|w|\sqrt{\hat\jmath(1-\hat\jmath)}$ ($0.22$--$0.26$), whose inclusion probabilities saturate at $\pi_d=1$ as $b$ grows so that their variance falls faster than $1/b$. For the judge arms the predicted and realised total savings correlate at $0.83$ with a mean absolute error of $0.06$; the $\lambda$-fitted control variate is predicted within $0.05$ everywhere but DBpedia-Entity at $\eps=0.01$ (predicted $0.19$--$0.21$, realised $0.07$--$0.10$), where the reference arm's certification curve rises slowly and $\Jfifty$ is sensitive to its shape rather than its scale. Applied to the menu-allocation runs with the realised design objective $\max_j V_j/s_j^2$ in place of the variance, the same model fails (correlation $0.53$ over 29 cells): the oracle's objective is $1.5$--$2.2\times$ better than static sharing at the budget where static sharing reaches 50\% $\ACT$ ($13\times$ on CAsT at $\eps=0.01$), yet its post-pilot cost is $0.79$--$1.00$ of static sharing, and the plug-in design's objective is $2$--$14\times$ \emph{worse} in six of seven cells while its cost is $0.83$--$1.24$. The objective, averaged over draws, is dominated by the draws whose slack is smallest; those draws do not certify under any design at these budgets, so improving them does not move the median draw that sets $\Jfifty$.

\paragraph{Pilot-only predictions.} \texttt{81\_sampling\_baselines.py --dump\_draws} records, for every draw and arm, the unique-label cost, the certificate outcome, the pilot-only within-query variance of the arm's estimator at $b=4$ documents per query (the design evaluated on the pilot queries, whose human and judge labels are known), and the pilot's slack estimates; \texttt{94\_j50\_ci.py} bootstraps draw indices, the same for every arm, for the paired intervals of Section~\ref{sec:table2}; \texttt{97\_pilot\_rule.py} scores the pilot predictions against the realised $\Jfifty$. On the development collections the pilot-predicted variance ratio correlates $0.93$ with the realised post-pilot cost ratio over all arms at $\eps=0.02$ ($0.95$ over the judge arms); the decision ``saving above 5\%'' is right in 15 of 18 judge cells at $\eps=0.02$ and 14 at $\eps=0.01$, and the predicted post-pilot labels of the reference arm have interquartile ranges over pilots of a factor of two to three, which bounds how precisely the total saving can be foreseen from 20 queries.



\section{The held-out protocol and its record}\label{app:heldout}
The rule of Section~\ref{sec:table2} has three ingredients fixed on the four development collections: the reference arm (decision-weight sampling), the threshold ($5\%$ of the total cost) and the vote (a judge arm is recommended when the majority of pilots predict a saving above the threshold). Its inputs are pilot quantities only: the within-query variance of every design at $b=4$ documents per query on the pilot's labelled documents, the pilot's slack of each comparison, and the pilot's between-query variance of the paired difference. \texttt{81\_sampling\_baselines.py --predict\_only} writes these for every draw without auditing, replaying the random stream so that the later audit uses the same draws; \texttt{98\_heldout\_predict.py} turns them into the predictions that are committed; \texttt{97\_pilot\_rule.py} scores them after the audit. The post-pilot labels of the reference arm are predicted in four versions (\texttt{lib/costpredict.py}): v1 is $z^2 b\,v_w/\hat s^2$, the formula fixed on the development collections; v2 replaces it by the smallest $L$ at which the two-stage bound of the population estimand---known pilot part, without-replacement sample of the rest, between-query variance from the pilot---falls below the slack; v3 takes the maximum of v2 over the comparisons of the menu; v4 uses a cross-fitted slack (candidate chosen on one half of the pilot, slack measured on the other). For Qwen3-8B all four were committed before the audit (repository commit \texttt{694e00b}); for the reranker only v1 was (\texttt{54deba0}), and its first lock carried a formula error that affected the label predictions but not the cost ratios, which the record (\texttt{05\_results/heldout/LOCK\_NOTE.md}) discloses. The scores, with every cell uncensored by a budget of one query-equivalent, are in \texttt{PILOT\_RULE\_*.csv}: at pilot 20 the cost-ratio correlation over all arms is $0.96$/$0.97$ (reranker, $\eps=0.02$/$0.01$) and $0.94$/$0.95$ (Qwen3-8B); the decisions of v3 are right in 8 of 8 cells for each judge, v2 and v4 in 8 of 8 for the reranker and 4 of 8 for Qwen3-8B (Touch\'e, predicted $0.044$--$0.055$ against realised $0.027$--$0.040$, on either side of the threshold), v1 in 4 of 8 for each. The order of events matters for reading these numbers: the reranker audit on both collections was run and scored first, v2--v4 were written in response to it, the Qwen3-8B judgments were then generated and all four predictors locked before the Qwen3-8B audit, and v3 was designated the primary predictor after that audit; \texttt{97\_pilot\_rule.py} also scores the executable single-pilot rule (Section~\ref{sec:table2}), whose per-draw choice uses only that draw's pilot. The NeuCLIRBench test followed a written lock (\texttt{03\_data/PROSPECTIVE\_LOCK\_v0.7}): the lock text was committed before the pool was built, the pool and the Qwen3-8B judgments were generated, the pilot-only predictions of the three judges were committed, and the audit was run afterwards; every prediction, score and criterion is in \texttt{05\_results/heldout}. On the development collections the corrected predictors do not remove the under-prediction of the post-pilot labels on CAsT and DBpedia-Entity ($636$ and $642$ against $1{,}152$ and $1{,}514$ with v3), whose policies are close: there the pilot's slack is optimistic because the candidate was chosen on the same 20 queries, and the cross-fitted slack of v4 removes the optimism (no false recommendation in 36 development cells) at the price of instability when the slack is near zero.

\begin{figure}[h]\centering
\includegraphics[width=0.7\linewidth]{figures/F5_act_vs_budget.png}
\caption{Cumulative $\ACT$ rate of the split certificate against the audit budget on DBpedia-Entity (500 repeats, $\eps=0.01$). Judges change the slope, not the validity: wrong certificates $\le2/500$ for every curve, including the inverted (adversarial) judge.}\label{fig:F5}
\end{figure}

\section{A calculator for the query-level map}\label{app:calculator}
To test whether Figure~\ref{fig:F6} follows from a handful of population quantities, \texttt{96\_ppi\_calculator.py} runs the certificate code of the repository (candidate by largest lower bound on the validation half, one-sided $t$ bounds, PPI++ with $\lambda$ cross-fitted on halves of the validation set and clipped to $[0,1]$, Bonferroni over the six ordered pairs) on synthetic audits: three policy utilities per query drawn from a Gaussian whose pairwise difference variances equal the population values ($\sigma$ from the human-only standard errors of the $\rho$ analysis), whose means equal the population gaps, and a judge whose paired differences have the population $\rho$ of each pair, attenuated proportionally to the map's noise levels. Each cell uses the same $T$, $N_{\mathrm{unlab}}$ and $\eps$ as the map and 1{,}000 audits. Figure~\ref{fig:F9} compares the simulated with the measured $\ACT$ ratios at $T=90$. On DBpedia-Entity (64 cells) the humans-only $\ACT$ rates correlate at $0.99$ (the Gaussian understates them at the smallest tolerances, by $0.09$ on average) and the ratios at $0.85$, with a mean absolute error of $0.06$; simulation and measurement agree on whether the ratio exceeds $1.2$ in 92\% of cells, the simulation being slightly conservative (7 cells above $1.2$ against 12). On TREC DL 2021--23 (32 cells) the mean absolute error is $0.03$ and the correlation $0.74$, and the simulation reproduces the ratios below one of the poor judge ($0.88$--$0.97$). Three variants isolate mechanisms. With the population-optimal $\lambda$ in place of the cross-fitted one, DBpedia-Entity is reproduced equally well (correlation $0.87$) but the TREC DL losses vanish (correlation $0.12$; every ratio $\ge0.99$): the loss of always-PPI with a poor judge is the cost of estimating $\lambda$ from 22-query folds. Removing the clip changes little on either collection. Dividing the cross-fitted $\lambda$ by $1+n/N$ (the finite-$N$ optimum) predicts ratios of $1.2$ at $N_{\mathrm{unlab}}=50$ on DBpedia-Entity where the measured ratio is $0.95$--$0.98$, so the finite-$N$ formula, which assumes the coefficient is known, is not the right correction when it is estimated. The practical reading is that the map can be computed before an audit from a pilot's estimates of $\rho$, $\sigma$ and the gaps together with $n$, $N_{\mathrm{unlab}}$ and $\eps$; how well pilot estimates of those quantities serve is the question of Section~\ref{sec:table2}'s pre-audit rule, which the released code evaluates.

\begin{figure}[h]\centering
\includegraphics[width=0.9\linewidth]{figures/F9_calculator.png}
\caption{Simulated against measured $\ACT$ ratio, $T=90$, every cell of Figure~\ref{fig:F6}; circles use the cross-fitted $\lambda$ of the implementation, triangles the population-optimal $\lambda$. Left DBpedia-Entity, right TREC DL 2021--23.}\label{fig:F9}
\end{figure}

\section{Reproducibility notes}\label{app:repro}
All numbers are regenerated from row-level result files by the scripts in the repository (the revised cost table by \texttt{86\_table2\_v2.py} from the \texttt{\_v2} result files; the cost model of Appendix~\ref{app:costmodel} by \texttt{95\_cost\_model.py}, its intervals and the pre-audit rule by \texttt{94\_j50\_ci.py} and \texttt{97\_pilot\_rule.py} from the per-draw records of \texttt{RUN\_TRACK\_C.sh}; the menu-allocation and exact-certificate appendices from \texttt{83\_menu\_allocation.py} and \texttt{63\_planner\_v2.py} with the options recorded in the design log; the pilot-size table by \texttt{88\_pilot\_sweep.py} from the \texttt{--n\_train} and \texttt{--fixed\_cand} runs of \texttt{83}). Each method uses an independent random stream derived from the repeat index and the method name. An earlier version of this paper reported that judge-uncertainty sampling ``harms''; that finding came from an implementation without a pilot-estimated residual model, and the faithful implementation reported in Table~\ref{tab:J50} does not harm, so the general claim was withdrawn and only the calibrated-judge rule is reported as harmful. Batched causal-LM inference with left padding must pass explicit position ids; without them the relevance judgments on one collection (TREC-COVID) collapsed to a constant, which we detected from a 49/50-query anomaly, fixed, and re-ran for every collection (label changes $\le0.5\%$ elsewhere; reranker scores unaffected, checked against batch-size-1 scoring). The BCa lower bound on $\rho$ used by the adoption rule covered its target at $0.88$--$0.92$ for pilots of 30 and 50 queries on DBpedia-Entity ($0.86$ at 20), against $0.75$--$0.84$ for the Fisher-$z$ bound on skewed pairs. The lock documents for Section~\ref{sec:prereg} were written before the target data were accessed (the only pre-lock access was the count of judged queries and pairs). The public repository's history, however, begins with a single release commit that already contains the lock documents together with the results, so the ordering is documented by the dated lock files and the development log rather than by the public commit sequence; a reader who needs independent evidence of the timing should treat the pre-registration as self-reported.
\end{document}
